% from ega1-7-fr.tex
\oldpage[I]{161}
de $A$ contenus dans $\mathfrak{j}_x$
(\hyperref[0.1.2.6-fr]{0, 1.2.6}), les idéaux premiers minimaux de $A_x$
correspondent aux idéaux premiers minimaux de $A$ contenus dans
$\mathfrak{j}_x$, d'où le lemme (\hyperref[I.1.1.14-fr]{1.1.14}).

Cela étant, supposons que l'on soit dans l'un des deux cas précités. Si $U$
est le domaine de définition de la $S$-application rationnelle $f$, désignons
par $f'$ l'application rationnelle de $\operatorname{Spec}(\mathscr{O}_x)$
dans $Y$ qui coïncide (compte tenu de (\hyperref[I.2.4.2-fr]{2.4.2})) avec
$f$ dans $U\cap\operatorname{Spec}(\mathscr{O}_x)$~; nous dirons que cette
application rationnelle est \emph{induite par $f$}.

\begin{proposition}[7.2.9]
\label{I.7.2.9-fr}
Soient $S$ un préschéma localement noethérien, $X$ un $S$-préschéma réduit,
$Y$ un $S$-schéma de type fini. On suppose en outre $X$ irréductible ou
localement noethérien. Soient alors $f$ une $S$-application rationnelle de
$X$ dans $Y$, $x$ un point de $X$. Pour que $f$ soit définie au point $x$,
il faut et il suffit que l'application rationnelle $f'$ de
$\operatorname{Spec}(\mathscr{O}_x)$ dans $Y$, induite par $f$
(\hyperref[I.7.2.8-fr]{7.2.8}), soit un morphisme.
\end{proposition}

La condition étant évidemment nécessaire (puisque
$\operatorname{Spec}(\mathscr{O}_x)$ est contenu dans tout ouvert contenant
$x$), prouvons qu'elle est suffisante. En vertu de
(\hyperref[I.6.5.1-fr]{6.5.1}), il existe un voisinage ouvert $U$ de $x$ dans
$X$ et un $S$-morphisme $g$ de $U$ dans $Y$, induisant $f'$ sur
$\operatorname{Spec}(\mathscr{O}_x)$. Si $X$ est irréductible, $U$ est dense
dans $X$, et en vertu de (\hyperref[I.7.2.3-fr]{7.2.3}) on peut supposer que
$g$ est une $S$-application rationnelle. En outre, le point générique de $X$
appartient à $\operatorname{Spec}(\mathscr{O}_x)$ et au domaine de définition
de $f$, donc $f$ et $g$ coïncident en ce point, et par suite dans un ensemble
ouvert non vide de $X$ (\hyperref[I.6.5.1-fr]{6.5.1}). Mais comme $f$ et $g$
sont des $S$-applications rationnelles, elles sont identiques
(\hyperref[I.7.2.3-fr]{7.2.3}), donc $f$ est définie en $x$.

Si maintenant on suppose $X$ localement noethérien, on peut supposer $U$
noethérien~; il n'y a alors qu'un nombre fini de composantes irréductibles
$X_i$ de $X$ contenant $x$ (\hyperref[I.7.2.8.1-fr]{7.2.8.1}), et on peut
supposer que ce sont les seules rencontrant $U$, en remplaçant au besoin $U$
par un ouvert plus petit (puisqu'il n'y a qu'un nombre fini de composantes
irréductibles de $X$ rencontrant $U$, $U$ étant noethérien). On voit alors
comme ci-dessus que $f$ et $g$ coïncident dans un ouvert non vide de chacun
des $X_i$. Tenant compte du fait que chacun des $X_i$ est contenu dans
$\overline U$, considérons alors le morphisme $f_1$, défini dans un ouvert
dense de $U\cup(X-\overline U)$, égal à $g$ dans $U$ et à $f$ dans
l'intersection de $X-\overline U$ et du domaine de définition de $f$. Comme
$U\cup(X-\overline U)$ est dense dans $X$, $f_1$ et $f$ coïncident dans un
ouvert dense de $X$, et comme $f$ est une application rationnelle, $f$ est
une extension de $f_1$ (\hyperref[I.7.2.3-fr]{7.2.3}), donc est définie au
point $x$.

\subsection{Faisceau des fonctions rationnelles.}
\label{subsection:I.7.3-fr}

\begin{env}[7.3.1]
\label{I.7.3.1-fr}
Soit $X$ un préschéma. Pour tout ouvert $U\subset X$, désignons par $R(U)$
l'anneau des fonctions rationnelles sur $U$
(\hyperref[I.7.1.3-fr]{7.1.3})~; c'est une
$\Gamma(U,\mathscr{O}_X)$-algèbre. En outre, si $V\subset U$ est un second
ouvert de $X$, toute section de $\mathscr{O}_X$ au-dessus d'une partie
ouverte partout dense de $U$ donne par restriction à $V$ une section
au-dessus d'une partie ouverte partout dense de $V$, et si deux sections
coïncident au-dessus d'une partie ouverte partout dense de $U$, leurs
restrictions à $V$ coïncident au-dessus d'une partie ouverte partout dense de
$V$. On définit donc ainsi un di-homomorphisme d'algèbres
$\rho_{V,U}:R(U)\to R(V)$, et il est clair que si $U\supset V\supset W$ sont
trois ouverts de $X$, on a
$\rho_{W,U}=\rho_{W,V}\circ\rho_{V,U}$~; les $R(U)$ définissent donc un
\emph{préfaisceau} d'algèbres sur $X$.
\end{env}

